\documentclass[11pt]{article}
\usepackage{amsmath, amssymb}
% \usepackage{natbib}
\usepackage{geometry}
\usepackage{float}
\usepackage[colorlinks=true, linkcolor=blue, citecolor=blue, urlcolor=blue]{hyperref}
\usepackage{setspace}
\usepackage{booktabs}
\usepackage{graphicx}
\usepackage{apacite}
\usepackage{natbib}
\usepackage{siunitx}
\usepackage{empheq}

\bibliographystyle{apacite}

\geometry{margin=1in}

\setstretch{1.2}
 \title{What happens to the exchange rate following a tariffs shock? Theory and some evidence}


\author{
Jason Lu\thanks{International Monetary Fund, Research Department. Email: {\tt jlu2@imf.org}} 
\and 
Dimitre Milkov\thanks{International Monetary Fund, Research Department. Email: {\tt dmilkov@imf.org}}
}

\date{First draft: November 10, 2025 \\ This version: \today}

\begin{document}
\maketitle

\begin{abstract}

An important feature of standard two-country trade models is the broad theoretical implication that a permanent increase in tariffs lead exchange rate appreciation for the country imposing the tariffs vs the tariffed country. However, we show that in a three-country setup, the effect on the exchange rate is ambiguous. 

\end{abstract}



\maketitle

\section{Introduction}



\section*{Exchange Rate Effects of a Tariff in a Two-Country Model}

\subsection*{1. Basic setup}

Let countries \( A \) and \( B \) trade goods.  
Define:
\[
E = \text{units of A's currency per unit of B's currency},
\]
so that an increase in \( E \) means an appreciation of B's currency.

Let \( P_A, P_B \) denote the domestic price levels in each country.  
Then the price of B's goods in A's currency is:
\[
P_B^* = E P_B.
\]

Exports from A to B and imports into A from B are given by:
\[
X_{AB} = X_{AB}(E), \qquad M_{AB} = M_{AB}(E),
\]
where \( X_{AB}'(E) > 0 \) and \( M_{AB}'(E) < 0 \).

Net exports of A are:
\[
NX_A(E) = X_{AB}(E) - M_{AB}(E),
\]
and by definition,
\[
NX_A(E) = -NX_B(E).
\]

\subsection*{2. Foreign exchange market equilibrium}

In equilibrium, demand and supply for B's currency must be equal:
\[
M_{AB}(E) = X_{AB}(E).
\]

\subsection*{3. Introducing a tariff}

Suppose country A imposes an ad valorem tariff \( \tau \) on imports from B.  
The domestic-currency price of B's goods in A becomes:
\[
P_B^T = (1+\tau) E P_B.
\]

This increases the effective price of imports and reduces import demand:
\[
M_{AB}^T(E) = M_{AB}\big( (1+\tau) E \big), \quad
M_{AB}^T(E) < M_{AB}(E).
\]

\subsection*{4. Exchange rate adjustment}

In the foreign exchange market:
\begin{align*}
\text{Demand for B's currency} &= M_{AB}^T(E), \\
\text{Supply of B's currency} &= X_{AB}(E).
\end{align*}

The tariff shifts the demand curve for B's currency leftward:
\[
M_{AB}^T(E) < M_{AB}(E).
\]

Equilibrium requires:
\[
M_{AB}^T(E^T) = X_{AB}(E^T),
\]
and since demand falls, the new equilibrium exchange rate satisfies:
\[
E^T < E^*,
\]
that is, B's currency \emph{depreciates} and A's currency \emph{appreciates}.

\subsection*{5. Summary}

\begin{center}
\begin{tabular}{lcc}
\hline
Country & Behavioral change & Currency effect \\
\hline
A (tariff-imposing) & $\downarrow$ import demand, $\uparrow$ net exports & Appreciates \\
B (exporting) & $\downarrow$ export revenue & Depreciates \\
\hline
\end{tabular}
\end{center}



\section{Setup}

Consider three countries indexed by \( i \in \{1,2,3\} \).  
Each country produces two types of goods:
\begin{itemize}
    \item A tradable good \( T_i \), differentiated by origin,
    \item A non-tradable good \( N_i \), consumed only domestically.
\end{itemize}

Labor is the only factor of production.

\section{Production}

Each good is produced with linear technology:
\begin{align}
Y_{T_i} &= A_{T_i} L_{T_i}, \\
Y_{N_i} &= A_{N_i} L_{N_i},
\end{align}
where \( A_{T_i} \) and \( A_{N_i} \) are sectoral productivities.  
The labor market clears in each country:
\begin{equation}
L_{T_i} + L_{N_i} = L_i.
\end{equation}

\section{Prices and Wages}

Let:
\begin{itemize}
    \item \( P_{T_i} \): world price of tradable good \( i \),
    \item \( P_{N_i} \): domestic price of non-tradables,
    \item \( w_i \): wage in country \( i \).
\end{itemize}

Under perfect competition, wages equal the marginal value product of labor:
\begin{equation}
w_i = P_{T_i} A_{T_i} = P_{N_i} A_{N_i}.
\end{equation}

Hence, the relative price of non-tradables (the real exchange rate) satisfies:
\begin{equation}
\frac{P_{N_i}}{P_{T_i}} = \frac{A_{T_i}}{A_{N_i}}.
\end{equation}

\section{Preferences}

Consumers in each country \( i \) have Cobb--Douglas preferences over all tradable goods and their own non-tradable good:
\begin{equation}
U_i = \left( \prod_{j=1}^3 C_{T_{ij}}^{\alpha_j} \right)
C_{N_i}^{1 - \sum_{j=1}^3 \alpha_j},
\end{equation}
where \( C_{T_{ij}} \) denotes consumption in country \( i \) of tradable good from country \( j \), and \( \alpha_j \in (0,1) \) is the expenditure share on tradable good \( j \).

\section{Budget Constraint}

Household income equals wage earnings:
\begin{equation}
w_i L_i = \sum_{j=1}^3 P_{T_j} C_{T_{ij}} + P_{N_i} C_{N_i}.
\end{equation}

\section{Market Clearing Conditions}

\subsection*{Labor Market}
\begin{equation}
L_{T_i} + L_{N_i} = L_i.
\end{equation}

\subsection*{Non-Tradables (Domestic Market Clearing)}
\begin{equation}
Y_{N_i} = C_{N_i}.
\end{equation}

\subsection*{Tradables (World Market Clearing)}
For each tradable good \( j \):
\begin{equation}
Y_{T_j} = \sum_{i=1}^3 C_{T_{ij}}.
\end{equation}

\section{Trade Balance}

The trade balance of country \( i \) is given by:
\begin{equation}
NX_i = \sum_{j=1}^3 P_{T_i} C_{T_{ji}} - \sum_{j=1}^3 P_{T_j} C_{T_{ij}}.
\end{equation}

Under balanced trade,
\begin{equation}
NX_i = 0 \quad \forall i.
\end{equation}

\section{Relative Prices and Real Exchange Rates}

Using the wage equalities:
\begin{equation}
\frac{w_i}{w_k} = \frac{P_{T_i} A_{T_i}}{P_{T_k} A_{T_k}},
\end{equation}
and
\begin{equation}
\frac{P_{N_i}}{P_{T_i}} = \frac{A_{T_i}}{A_{N_i}},
\end{equation}
we see that countries with higher tradable productivity \( A_{T_i} \) have more appreciated real exchange rates (a higher \( P_{N_i}/P_{T_i} \)).

\section{Summary of Equilibrium Conditions}

An equilibrium consists of prices \( \{P_{T_i}, P_{N_i}\} \), wages \( \{w_i\} \), and allocations \( \{L_{T_i}, L_{N_i}, C_{T_{ij}}, C_{N_i}\} \) such that:
\begin{enumerate}
    \item Firms maximize profits: \( w_i = P_{T_i} A_{T_i} = P_{N_i} A_{N_i} \).
    \item Households maximize utility given their budget constraint.
    \item Labor markets clear: \( L_{T_i} + L_{N_i} = L_i \).
    \item Goods markets clear: \( Y_{N_i} = C_{N_i} \) and \( Y_{T_j} = \sum_i C_{T_{ij}} \).
    \item Optional: \( NX_i = 0 \) for all \( i \).
\end{enumerate}

\section{Extensions}

Possible extensions include:
\begin{itemize}
    \item Introducing trade costs or tariffs between countries.
    \item Allowing capital accumulation or investment.
    \item Incorporating intermediate inputs or sectoral linkages.
    \item Adding intertemporal dynamics and international borrowing/lending.
\end{itemize}

\documentclass[12pt]{article}
\usepackage{amsmath, amssymb, geometry}
\geometry{margin=1in}

\begin{document}

\title{A Three-Country Trade Model with Tradables, Non-Tradables, and Labor}
\author{}
\date{}
\maketitle

\section{Setup}

We consider three countries \( i \in \{1,2,3\} \), each producing:
\begin{itemize}
    \item A \textit{tradable good} \( T_i \), differentiated by origin.
    \item A \textit{non-tradable good} \( N_i \), consumed only domestically.
\end{itemize}

Labor is the only factor of production, immobile across countries.

\section{Production}

Each country has linear production technologies:
\begin{align}
Y_{T_i} &= A_{T_i} L_{T_i}, \\
Y_{N_i} &= A_{N_i} L_{N_i},
\end{align}
where \( A_{T_i} \) and \( A_{N_i} \) denote productivities in the tradable and non-tradable sectors, respectively.  
Labor market clearing implies:
\begin{equation}
L_{T_i} + L_{N_i} = L_i.
\end{equation}

\section{Prices and Wages}

Let:
\begin{itemize}
    \item \( P_{T_i} \): price of tradable good \( i \) on the world market,
    \item \( P_{N_i} \): domestic price of non-tradables,
    \item \( w_i \): wage in country \( i \).
\end{itemize}

Under perfect competition:
\begin{equation}
w_i = P_{T_i} A_{T_i} = P_{N_i} A_{N_i}.
\end{equation}

Hence, the relative price of non-tradables is:
\begin{equation}
\frac{P_{N_i}}{P_{T_i}} = \frac{A_{T_i}}{A_{N_i}}.
\label{eq:real_exchange_rate}
\end{equation}
Equation \eqref{eq:real_exchange_rate} shows that countries with higher tradable productivity have more appreciated real exchange rates.

\section{Preferences and Demand}

Consumers in each country \( i \) have Cobb--Douglas preferences:
\begin{equation}
U_i = \left( \prod_{j=1}^3 C_{T_{ij}}^{\alpha_j} \right)
C_{N_i}^{1 - \sum_{j=1}^3 \alpha_j},
\end{equation}
where \( C_{T_{ij}} \) is consumption in country \( i \) of tradable good \( j \), and \( \alpha_j \in (0,1) \) is the expenditure share on good \( j \).

Given income \( w_i L_i \), the consumer’s budget constraint is:
\begin{equation}
w_i L_i = \sum_{j=1}^3 P_{T_j} C_{T_{ij}} + P_{N_i} C_{N_i}.
\end{equation}

Cobb--Douglas preferences imply the following expenditure shares:
\begin{align}
P_{T_j} C_{T_{ij}} &= \alpha_j w_i L_i, \quad \forall j, \\
P_{N_i} C_{N_i} &= \left( 1 - \sum_{j=1}^3 \alpha_j \right) w_i L_i.
\end{align}

\section{Market Clearing}

\subsection*{Tradable Goods}

Each tradable good \( j \) is sold worldwide:
\begin{equation}
Y_{T_j} = \sum_{i=1}^3 C_{T_{ij}}.
\label{eq:tradables_clearing}
\end{equation}

Substitute the demand function:
\begin{equation}
A_{T_j} L_{T_j} = \sum_{i=1}^3 \frac{\alpha_j w_i L_i}{P_{T_j}}.
\end{equation}

Rearranging gives the world market equilibrium condition for the price of tradable good \( j \):
\begin{equation}
P_{T_j} = \frac{\sum_{i=1}^3 \alpha_j w_i L_i}{A_{T_j} L_{T_j}}.
\label{eq:tradable_price}
\end{equation}

\subsection*{Non-Tradables}

Non-tradables must clear domestically:
\begin{equation}
A_{N_i} L_{N_i} = C_{N_i} = \frac{(1 - \sum_j \alpha_j) w_i L_i}{P_{N_i}}.
\end{equation}

Substituting \( P_{N_i} = \frac{w_i}{A_{N_i}} \) yields:
\begin{equation}
A_{N_i} L_{N_i} = (1 - \sum_j \alpha_j) A_{N_i} L_i,
\end{equation}
which implies that the share of labor employed in non-tradables equals the expenditure share:
\begin{equation}
\frac{L_{N_i}}{L_i} = 1 - \sum_j \alpha_j, \quad
\frac{L_{T_i}}{L_i} = \sum_j \alpha_j.
\end{equation}

\section{Equilibrium Wages and Relative Prices}

Using \( w_i = P_{T_i} A_{T_i} \) and substituting \eqref{eq:tradable_price}:
\begin{align}
w_i &= A_{T_i} P_{T_i} = A_{T_i} 
\frac{\sum_{k=1}^3 \alpha_i w_k L_k}{A_{T_i} L_{T_i}} 
= \frac{\sum_{k=1}^3 \alpha_i w_k L_k}{L_{T_i}}.
\end{align}

Assuming balanced trade and symmetric preferences (i.e.\ \( \alpha_j = \alpha/3 \) for all \( j \)), and symmetric tradable labor allocation \( L_{T_i} = \alpha L_i \), we obtain:
\begin{equation}
w_i = \bar{w} = \text{constant across } i,
\end{equation}
so that tradable prices are proportional to productivity:
\begin{equation}
P_{T_i} = \frac{\bar{w}}{A_{T_i}}.
\end{equation}

Finally, using \eqref{eq:real_exchange_rate}, the relative price of non-tradables becomes:
\begin{equation}
\frac{P_{N_i}}{P_{T_i}} = \frac{A_{T_i}}{A_{N_i}}
\quad \Rightarrow \quad
P_{N_i} = \frac{\bar{w}}{A_{N_i}}.
\end{equation}

\section{Interpretation}

\begin{itemize}
    \item Countries with higher tradable productivity \( A_{T_i} \) have lower \( P_{T_i} \) (their tradables are cheaper internationally) and higher real exchange rates \( P_{N_i}/P_{T_i} \).
    \item The real exchange rate depends on relative productivity across sectors: a Balassa--Samuelson effect.
    \item Under symmetry and balanced trade, wages equalize internationally.
\end{itemize}

\section{Summary of Equilibrium Relationships}

\begin{align}
w_i &= P_{T_i} A_{T_i} = P_{N_i} A_{N_i}, \\
\frac{P_{N_i}}{P_{T_i}} &= \frac{A_{T_i}}{A_{N_i}}, \\
P_{T_i} &= \frac{\bar{w}}{A_{T_i}}, \quad P_{N_i} = \frac{\bar{w}}{A_{N_i}}, \\
\frac{L_{N_i}}{L_i} &= 1 - \sum_j \alpha_j.
\end{align}

\end{document}



\end{document}


\end{document}
